\chapter{Conclusion}
Moderation analysis has been a core statistical tool within psychology for a number of years; it is certainly recognized as a staple within our methodological toolkit. Psychological networks represent a much newer class of methods, although have gained incredible popularity and widespread use in the short time since they have been developed. Moderated networks offer an integration of these two frameworks in a way that is both incremental (e.g., adding interaction terms to nodewise regression models; treating moderators as exogenous variables) and novel in its approach (e.g., visualizing conditional networks; performing global model comparison and goodness-of-fit tests). Moreover, while the \code{modnets} package for R is still being developed, it already contains many different estimation procedures and tools for both plotting and analyzing results from various types of MNMs. 
